%% =====================================================================
%%  Bilan 2 — La vie de tous les jours (unités 5 à 8)
%%  Pistes : B2.01 (écoute notée) · B2.02 (expression orale)
%% =====================================================================
\bilan{2}

\noindent{\fontsize{11}{14.6}\selectfont Deuxième partie terminée ! Vous savez maintenant décrire
des objets, raconter votre journée, dire ce que vous aimez et présenter votre logement. Révisez
l'essentiel, puis faites les huit exercices notés sur 20. Les corrigés sont
page~\pageref{cor-b2-1}.\par}

%% ---------------------------------------------------------------------
\rubrique{grammaire}{Révisez}

\begin{paire}
\begin{regle}[Les noms et les articles · unité 5]
\en{a pen} / \en{an apple} · \en{the} = le, la, les (chose précise)\par\vspace{1mm}
Pluriel : \en{pens}, \en{boxes}, \en{cities}, \en{men}, \en{women}, \en{children}, \en{people}\par\vspace{1mm}
\en{this} / \en{these} (ici) · \en{that} / \en{those} (là-bas)
\end{regle}
\begin{regle}[Le présent simple · unités 6 et 7]
\enb{I work} · \enb{she work\textbf{s}}\par
\enb{I don't work} · \enb{she doesn't work}\par
\enb{Do you work?} · \enb{Does she work?}\par
\enb{Yes, I do.} · \enb{No, she doesn't.}\par\vspace{1mm}
{\fontsize{9.6}{12}\selectfont Un seul \en{-s} par phrase : \en{she works}, mais \en{she doesn't work}.\par}
\end{regle}
\end{paire}

\begin{paire}
\begin{regle}[Dire l'heure et les habitudes · unité 6]
\en{What time is it? — It's nine thirty.}\par\vspace{1mm}
\en{at seven} · \en{in the morning} · \en{on Mondays} · \en{every day}\par\vspace{1mm}
Goûts : \en{like}, \en{love}, \en{hate} + nom ou \en{-ing}
\end{regle}
\begin{regle}[\en{there is} et les prépositions · unité 8]
\en{There's a kitchen.} · \en{There are two windows.}\par
\en{Is there…? — Yes, there is.}\par\vspace{1mm}
\en{in} · \en{on} · \en{under} · \en{next to} · \en{behind} · \en{in front of} · \en{between} · \en{opposite}
\end{regle}
\end{paire}

%% ---------------------------------------------------------------------
\rubrique{pratique}{Testez-vous}

\exercice[ecrire,niveau=1,points=2]{Le pluriel}{Écrivez le pluriel. \textit{(0,5 point par réponse)}}
\begin{questions}[4]
\item box → \trou[22mm]{boxes}
\item city → \trou[22mm]{cities}
\item child → \trou[22mm]{children}
\item woman → \trou[22mm]{women}
\end{questions}

\exercice[ecrire,niveau=2,points=3]{La journée de Mrs.\ Lopez}{Complétez avec les verbes au présent
simple. \textit{(0,5 point par réponse)}}
\begin{anglais}
\begin{tcolorbox}[otBase,colback=otNavyTint,colframe=otNavyTint,fontupper=\fontsize{11}{19}\selectfont]
Mrs.\ Lopez \trou[22mm]{gets up} \textit{(get up)} at seven. She \trou[22mm]{has} \textit{(have)} breakfast
with Miguel. Then she \trou[22mm]{goes} \textit{(go)} to the market. In the evening, she
\trou[22mm]{watches} \textit{(watch)} TV. She \trou[28mm]{doesn't work} \textit{(not work)} on Sundays,
and she \trou[22mm]{loves} \textit{(love)} cooking for her family.
\end{tcolorbox}
\end{anglais}

\exercice[ecrire,niveau=2,points=3]{Questions et réponses}{Écrivez la question (1 à 3) ou la réponse
courte (4 à 6). \textit{(0,5 point par réponse)}}
\begin{questions}
\item \textit{(you / like tea)} \hfill \trou[70mm]{Do you like tea?}
\item \textit{(Emma / drink coffee)} \hfill \trou[70mm]{Does Emma drink coffee?}
\item \textit{(they / live in Brooklyn)} \hfill \trou[70mm]{Do they live in Brooklyn?}
\item Does Jack cook? — \hfill \trou[50mm]{No, he doesn't.}
\item Do you like jazz? — \hfill \trou[50mm]{Yes, I do.}
\item Is there a garden? — \hfill \trou[50mm]{No, there isn't.}
\end{questions}

\exercice[lire,ecrire,niveau=1,points=2]{Où est la balle ?}{Complétez avec la bonne préposition.
\textit{(0,5 point par réponse)}}
\begin{tcbraster}[raster columns=4,raster equal height=rows,raster column skip=2.4mm,raster before skip=1mm]
\otCartePrepo{under}{It's}{under}{the table.}
\otCartePrepo{in}{It's}{in}{the box.}
\otCartePrepo{between}{It's}{between}{the boxes.}
\otCartePrepo{behind}{It's}{behind}{the box.}
\end{tcbraster}

\exercice[lire,ecrire,niveau=1,points=2]{L'heure}{Écrivez l'heure en anglais.
\textit{(0,5 point par réponse)}}
\begin{tcbraster}[raster columns=4,raster equal height=rows,raster column skip=2.4mm,raster before skip=1mm]
\otCarteHeure{7}{30}{7 h 30}{}{soleil}{It's seven thirty.}
\otCarteHeure{11}{15}{11 h 15}{}{soleil}{It's eleven fifteen.}
\otCarteHeure{4}{45}{16 h 45}{}{soleil}{It's four forty-five.}
\otCarteHeure{9}{0}{21 h 00}{}{lune}{It's nine o'clock.}
\end{tcbraster}

\exercice[ecouter,niveau=2,points=3,piste=B2.01]{Chez Sam}{Sam présente son appartement et sa
journée. Écoutez, puis indiquez si chaque phrase est vraie (V) ou fausse (F).
\textit{(0,5 point par réponse)}}
\begin{questions}[2]
\item There are two bedrooms. \hfill\trou[9mm]{F}
\item There's a balcony. \hfill\trou[9mm]{F}
\item The guitars are in the living room. \hfill\trou[9mm]{V}
\item Sam gets up at five. \hfill\trou[9mm]{V}
\item Maya lives with Sam. \hfill\trou[9mm]{V}
\item Sam sleeps in the bedroom. \hfill\trou[9mm]{F}
\end{questions}

\exercice[lire,niveau=2,points=3]{Le blog de Léa}{Lisez le message de Léa, puis indiquez si chaque
phrase est vraie (V) ou fausse (F). \textit{(0,5 point par réponse)}}
\begin{anglais}
\begin{tcolorbox}[otBase,colback=white,colframe=otRule,boxrule=0.5pt,fontupper=\fontsize{10.8}{14.4}\selectfont]
{\otSemi\fontsize{11}{13}\selectfont My life in New York}\par\vspace{1.4mm}
I get up at seven thirty every day. I don't take the subway: I walk to the studio. It's twenty
minutes. I start work at nine. At lunch, Jack and I eat sandwiches, but Emma doesn't like
sandwiches: she has a salad. In the evening, I cook or I go to Sam's coffee shop. On weekends, I
love walking in Central Park. I don't like the subway on Saturdays: there are too many people!
\par\vspace{1mm}{\fontsize{8.6}{10}\selectfont\color{otMuted}\foreignlanguage{french}{\textit{minutes} : minutes · \textit{a salad} : une salade · \textit{too many people} : trop de monde}}
\end{tcolorbox}
\begin{questions}[2]
\item Léa takes the subway to work. \hfill\trou[9mm]{F}
\item Léa starts work at nine. \hfill\trou[9mm]{V}
\item Emma eats sandwiches at lunch. \hfill\trou[9mm]{F}
\item Léa likes walking in Central Park. \hfill\trou[9mm]{V}
\item Léa likes the subway on Saturdays. \hfill\trou[9mm]{F}
\item The studio is twenty minutes from Léa's apartment. \hfill\trou[9mm]{V}
\end{questions}
\end{anglais}

\exercice[ecrire,niveau=3,points=2]{Votre dimanche}{En cinq phrases, décrivez votre logement et ce que
vous faites le dimanche. \textit{(0,5 point par phrase correcte, 2 points au maximum)}}
\lignes{5}

%% ---------------------------------------------------------------------
\rubrique{ecouter}{Parlez}

\exercice[ecouter,parler,niveau=3,piste=B2.02]{Une journée ordinaire}{Emma vous pose des questions
sur votre vie quotidienne. Répondez-lui \textbf{à voix haute} pendant chaque pause. Cet exercice n'est
pas noté : comparez vos réponses avec les modèles du corrigé.}
\begin{center}\otTableau
\begin{tabular}{@{}>{\otSemi\fontsize{8.6}{10}\selectfont\color{otNavy}}l@{\hspace{3mm}}>{\enOn}p{70mm}<{\enOff}@{\hspace{6mm}}p{78mm}@{}}
 & \otTH{Emma dit\ldots} & \otTH{Vous\ldots}\\\midrule
1. & What time do you get up? & donnez l'heure de votre lever.\\
2. & Do you take the subway or the bus? & répondez.\\
3. & What do you have for breakfast? & dites ce que vous mangez et buvez.\\
4. & Is there a coffee shop near your home? & répondez par une réponse courte.\\
5. & What do you do on Sundays? & décrivez votre dimanche en une ou deux phrases.\\\bottomrule
\end{tabular}\end{center}

\otScore{20}
